\section{Discussion}
\label{sec:limitations}
%This needs to go somewhere for the checklist.
%\subsection{Principle findings}

We performed a systematic review of 445 benchmarks from the NLP and ML literature to assess the best practices around construct validity. We found that the operationalisation of abstract phenomena was often insufficient, with definitions being missing or contested. Tasks were frequently taken from pre-existing data sources without adjustments to ensure that they were representative of the target phenomenon. Statistical testing was also rarely performed. About half of the reviewed articles did discuss the validity of their benchmark, but nearly every paper had weaknesses in at least one area. In light of these gaps, we created a list of recommendations covering the design of phenomena, tasks, and metrics as well as interpretation to improve the construct validity of future LLM benchmarks.

%\subsection{How to use the Checklist}
We developed the operational checklist as a practical tool to support researchers in proactively engaging with construct validity throughout the benchmark lifecycle. We recommend its use early in the design phase to guide critical design decisions regarding task selection, sample construction, metric justification, as well as later when interpreting findings. We do not expect that every benchmark will satisfy every item, as practical trade-offs will sometimes be necessary. We recommend reporting the checklist as an appendix with answers and explanations for each skipped item, allowing users to assess if a benchmark aligns with their needs and matches best practices.
Beyond new benchmark developments, the checklist can also serve as an evaluation framework for existing benchmarks, or for adapting them to new domains or capabilities.

%\subsection{Limitations}\label{sec:limitations}

% \paragraph{Selection Bias} 
We describe limitations of our approach. Our focus on leading conference proceedings, while ensuring a baseline of peer-reviewed quality, may systematically exclude certain types of impactful benchmarks. For example, it does not capture benchmarks developed and released by industry labs without formal peer review, or those published in specialised domain-specific venues, may not be captured. We primarily review benchmarks prevalent in mainstream academic AI research.

% \paragraph{LLM Filtering} 
To manage the extensive initial corpus, we employed GPT-4o mini for preliminary screening against topic, empirical, and modality criteria, prior to manual review. This automated step was only used to exclude articles, and validated against human annotation (see~\cref{app:schema} indicating good agreement). Nevertheless, this may have introduced undetected false negative systematic errors. Distributional shift in language usage may also have contributed to the lower inclusion of older papers, alongside the general increase in papers being published in this area. The scale of the review also necessitated limiting the number of reviewers per paper, reducing the robustness of the reviews.

\section{Conclusion}
The rapid advancement of LLMs requires robust evaluation. Our systematic review of 445 benchmarks reveals prevalent gaps that undermine the construct validity needed to accurately measure targeted phenomena. To address these shortcomings, which can hinder genuine progress, we propose eight recommendations and a practical checklist for designing and interpreting LLM benchmarks. Ultimately, ``measuring what matters'' requires a conscious, sustained effort from the research community to prioritise construct validity, fostering a cultural shift towards more explicit and rigorous validation of evaluation methodologies.
